\DocumentMetadata{
  pdfversion=2.0,
  pdfstandard=ua-2,
  lang=en-US,
  tagging=on,
  tagging-setup={math/setup=mathml-SE,tabsorder=structure}
}
\documentclass[11pt,letterpaper]{article}
\usepackage[margin=0.68in,includefoot]{geometry}
\usepackage{newcomputermodern}
\usepackage{amsmath,amscd,mathtools}
\usepackage{tikz,adjustbox}
\providecommand{\square}{\mdlgwhtsquare}
\usepackage[hidelinks]{hyperref}
\hypersetup{
  pdftitle={Numerical Analysis PhD Exam, August 2025},
  pdfauthor={University of Florida Department of Mathematics},
  pdfsubject={Graduate mathematics examination},
  pdfdisplaydoctitle=true,
  bookmarksopen=true
}
\setcounter{secnumdepth}{0}
\setlength{\parindent}{0pt}
\setlength{\parskip}{0.28em}
\setlength{\emergencystretch}{3em}
\setlength{\leftmargini}{2.15em}
\setlength{\leftmarginii}{2.65em}
\setlength{\leftmarginiii}{3em}
\renewcommand{\labelenumi}{\textbf{\theenumi.}}
\pagestyle{empty}
\raggedbottom
\allowdisplaybreaks
\newenvironment{examproblems}[1]{%
  \begin{enumerate}%
  \setcounter{enumi}{#1}%
  \setlength{\topsep}{0.35em}%
  \setlength{\partopsep}{0pt}%
  \setlength{\itemsep}{0.48em}%
  \setlength{\parsep}{0.18em}%
}{\end{enumerate}}
\newenvironment{examsubparts}{%
  \begin{enumerate}%
  \setlength{\topsep}{0.18em}%
  \setlength{\partopsep}{0pt}%
  \setlength{\itemsep}{0.16em}%
  \setlength{\parsep}{0.1em}%
}{\end{enumerate}}
\newenvironment{examinstructions}{%
  \par\begingroup\itshape
}{%
  \par\endgroup\vspace{0.18em}
}
\makeatletter
\renewcommand\section{\@startsection{section}{1}{0pt}%
  {-1.25ex plus -.25ex minus -.1ex}%
  {0.55ex plus .1ex}%
  {\normalfont\large\bfseries}}
\renewcommand{\@maketitle}{%
  \newpage\null\vskip -1em%
  {\footnotesize\bfseries
    \href{https://www.ufl.edu/}{University of Florida}, \href{https://math.ufl.edu/}{Department of Mathematics}\par}%
  \vskip -0.5em%
  \tagstructbegin{tag=Title}%
    {\LARGE\bfseries\href{https://gma.math.ufl.edu/exams/numerical-analysis-phd/}{Numerical Analysis PhD Exam}, August 2025\par}%
  \tagstructend%
  \vskip 0.25em%
  \tagmcbegin{artifact}%
    \hrule height 1pt%
  \tagmcend%
  \vskip 1em%
}
\makeatother
\title{Numerical Analysis PhD Exam, August 2025}
\author{}
\date{}
\begin{document}
\maketitle
\thispagestyle{empty}
\begin{examinstructions}
Do 4 (four) of the first 5 problems (1–5) and 4 (four) of the last 5 problems (6–10).
\end{examinstructions}
\begin{examproblems}{0}
\item
Suppose~\(A\) is Hermitian positive definite.
\begin{examsubparts}
\item[\textnormal{(a)}]
Prove that each principal submatrix of~\(A\) is Hermitian positive definite.
\item[\textnormal{(b)}]
Prove that an element of~\(A\) with largest magnitude lies on the diagonal.
\item[\textnormal{(c)}]
Prove that~\(A\) has a Cholesky decomposition.
\end{examsubparts}
\item
Assume \(A\in\mathbb{C}^{m\times m}\).
\begin{examsubparts}
\item[\textnormal{(a)}]
Show that~\(A\) has a Schur decomposition.
\item[\textnormal{(b)}]
If~\(A\) has a collection of~\(m\) linearly independent eigenvectors, show that~\(A\) is diagonalizable.
\end{examsubparts}
\item
Let \(A\in\mathbb{C}^{m\times n}\), with \(m\ge n\) and \(\operatorname{rank}(A)=p=n\ge3\). Let \(a_1,a_2,\ldots\) denote the columns of~\(A\).
\begin{examsubparts}
\item[\textnormal{(a)}]
Using the modified Gram–Schmidt process, write out expressions for \(q_1,q_2,q_3\), the first three columns of~\(Q\) in the QR decomposition of~\(A\).
\item[\textnormal{(b)}]
Show the vector \(q_3\) found in part (a) is orthogonal to both \(q_1\) and \(q_2\).
\end{examsubparts}
\item
Let \(A=U\Sigma V^*\) be the singular value decomposition of \(A\in\mathbb{C}^{m\times n}\). Let \(u_j\) denote column~\(j\) of~\(U\).
\begin{examsubparts}
\item[\textnormal{(a)}]
Suppose \(\operatorname{rank}(A)=p<n<m\). Show \(\{u_1,u_2,\ldots,u_p\}\) is a basis for \(\operatorname{Col}(A)\) and \(\{u_{p+1},u_{p+2},\ldots,u_m\}\) is a basis for \(\operatorname{Null}(A^*)\).
\item[\textnormal{(b)}]
Suppose~\(A\) is full rank and \(x\ne0\). Let \(\sigma_i\), \(i=1,\ldots,n\), be the singular values of~\(A\). Show 
\[
\sigma_1\ge \frac{\lVert Ax\rVert_2}{\lVert x\rVert_2}\ge\sigma_n>0.
\]
 If you want to use the property that \(\lVert A\rVert_2=\sigma_1\), then you must prove that.
\end{examsubparts}
\item
If \(q_1,\ldots,q_n\in\mathbb{C}^m\) is an orthonormal basis for the subspace \(V\subset\mathbb{C}^m\), prove that the orthogonal projector onto~\(V\) is \(QQ^*\), where~\(Q\) is the matrix whose columns are the \(q_j\).
\item
Consider the fixed point iteration method \(x_{k+1}=g(x_k)\), \(k=0,1,\ldots\), for solving the nonlinear equation \(f(x)=0\). Consider choosing an iteration function of the form 
\[
g(x)=x-af(x)-b(f(x))^2-c(f(x))^3,
\]
 where \(a,b,\) and~\(c\) are parameters to be determined. Assume~\(f\) is sufficiently differentiable and the corresponding iterations \(x_{k+1}=g(x_k)\) converge to a unique fixed point~\(z\). Find expressions for the parameters \(a,b,\) and~\(c\) in terms of functions of~\(z\) such that the iteration method is of fourth order.
\item
Consider the Lagrange polynomial approximation \(p(x)\) of a function \(f(x)\in C^2[a,b]\) at two points \(x_0,x_1\in[a,b]\). Determine \(\alpha(x)\) explicitly in the following formula 
\[
f(x)-p(x)=\frac{f^{(n+1)}(\alpha(x))\prod_{i=0}^{n}(x-x_i)}{(n+1)!}
\]
 for the case 
\[
f(x)=\frac1x,\qquad x_0=1,\qquad x_1=2.
\]
\item
Assume that \(f\in C^3[a,b]\) and \(x,x+h,x+2h\in[a,b]\). Determine constants \(c_1\) and \(c_2\) such that 
\[
\left|f'(x)-\frac1h\left[-\frac32f(x)+c_1f(x+h)+c_2f(x+2h)\right]\right|\le ch^2\max_{a\le x\le b}|f'''(x)|,
\]
 where \(c>0\) is a constant independent of~\(h\).
\item
Let \(Q_m(f)\) denote the~\(m\)-point Gaussian quadrature rule over the interval \([a,b]\) and with continuous weight function \(\rho(x)\ge0\), that is, 
\[
Q_m(f)=\sum_{i=1}^{m}\alpha_i f(x_i)\approx J(f)=\int_a^b\rho(x)f(x)\,dx.
\]
 Show that, if~\(a\) and~\(b\) are finite and~\(f\) is continuous, then \(Q_m(f)\to J(f)\) as \(m\to\infty\).
\item
Consider numerically solving the initial value problem 
\[
y'(t)=f(t,y),\quad 0<t\le t_f,\qquad \text{with }y(0)=\eta.
\]
 Assume~\(f\) is sufficiently differentiable and let~\(h\) denote the step size. Prove that the method 
\[
y_{n+3}+y_{n+2}-y_{n+1}-y_n=h(f_{n+3}+f_{n+2}+f_{n+1}+f_n)
\]
 is consistent but not convergent.
\end{examproblems}
\end{document}
